\documentclass[12pt,a4paper]{article}
\usepackage[top=1.5cm, bottom=3cm, left=2cm, right=2cm]{geometry}
\usepackage{amsmath,amsthm,amsfonts,amssymb,amscd,mathtools}
\usepackage{enumerate}
\usepackage{enumitem}
\usepackage{fancyhdr}
\usepackage{fancyvrb}
\usepackage{fullpage}
\usepackage{hyperref}
\usepackage{lastpage}
\usepackage{mathrsfs}
\setlength{\parindent}{0.0in}
\setlength{\parskip}{0.05in}

\pagestyle{fancyplain}
\headheight 20pt
\lhead{\textbf{COMP1046-CW1}}              
\rhead{Student ID: \textbf{20717902
}}
\cfoot{}
\rfoot{\small Page \thepage}
\headsep 1em

\begin{document}
%%% Please start every new question in a new page %%%
\subsection*{Question 1}
Given:\\
$\forall x \exists y(C(x) \to (D(y) \land L(y,x)))$\\
$\forall x \forall y(D(x) \land S(y) \to \neg L(x,y))$\\
To prove: $\forall x(C(x) \to \neg S(x))$
\begin{table}[h!]
\centering
\renewcommand{\arraystretch}{1.2} 
\begin{tabular}{|l | p{6cm} | p{4.5cm} |}
\hline
\textbf{Number} & \textbf{Step} & \textbf{Reason} \\
\hline
1 & $\forall x \exists y(C(x) \to (D(y) \land L(y,x)))$& Premise(Given) \\
\hline
2 & $\forall x \forall y(D(x) \land S(y) \to \neg L(x,y))$ & Premise(Given) \\
\hline
3 & $ \exists y(C(a) \to (D(y) \land L(y,a)))$&Universal Instantiation from 1, $a$ is an arbitrary element\\
\hline
4&$C(a) \to (D(b) \land L(b,a))$&Existential Instantiation from 3, for some $b$\\
\hline
5&$D(b) \land L(b,a)$&Modus Ponens from 4, assuming C(a) is true \\
\hline
6&$L(b,a)$&Simplification from 5\\
\hline
7&$D(b)$ &Simplification from 5 \\
\hline
8& $D(b) \land S(a) \to \neg L(b,a)$ &Universal Instantiation from 2,$x=b$,  $y=a$\\
\hline
9&$\neg (D(b) \land S(a))$ &Modus Tollens from 6, 7 and 8 \\
\hline
10&$\neg D(b) \lor \neg S(a)$ &De Morgan's law \\
\hline
11&$\neg S(a)$&Disjunctive Syllogism from 7 and 10\\
\hline
12& $C(a) \to \neg S(a)$ & Conditional Proof from 5 and 11\\
\hline
13&$\forall x(C(x) \to \neg S(x))$& Universal Generalization from 12 \\
\hline
\end{tabular}
\end{table}
\newpage
\subsection*{Question 2}
Solution: Let P(n) be the propostion that $ 7^n - 6n - 1 \equiv 0 \pmod{36}$, for $n \in \mathbb{Z},n \geq 1$. 
BASIS STEP:\\ P(1)is true because
\begin{align*}
7^1 - 6(1) - 1 &= 7 - 6 - 1 = 0 \\
0 &= 0 \pmod{36} \quad 
\end{align*}
\\
INDUCTIVE STEP:\\
For the inductive hypothesis, we assume that P($k$) is true for an integer $k$ that $k \in \mathbb{Z} $  and $k \geq 1$. That is, we assume that
\[
7^k - 6k - 1 \equiv 0 \pmod{36}
\] for a integer $k \in \mathbb{Z}$ and $k \geq 1$.\\
To carry out the inductive step using this assumption, we assume that P(k) is true, then we must show that P(k+1) is also true, which is
\[
7^{k+1} - 6(k + 1) - 1 \equiv 0 \pmod{36}
\]
Because
\begin{align*}
7^{k+1} - 6(k + 1) - 1 &= 7 \cdot 7^k - 6k - 6 - 1 \\
&= 7 \cdot 7^k - 6k - 7
\end{align*}

So, $7 \cdot 7^k - 6k - 7 = 36n$ for an existing $n \in \mathbb{Z}$.

Using the inductive hypothesis, we need to prove that $7^{k+1} - 6(k + 1) - 1 \equiv 0 \pmod{36}$:

\begin{align*}
7^{k+1} - 6(k + 1) - 1 & = 7(36n + 6k + 1) - 6k - 7 \\
& = 252n + 42k + 7 - 6k - 7 \\
& = 36(7n + k)
\end{align*}
Because we have completed the basis step and the inductive step, by mathematical induction, we know that P(n) is true for all $n \in \mathbb{Z},n \geq 1$. That is $ 7^n - 6n - 1 \equiv 0 \pmod{36}$, for $n \in \mathbb{Z},n \geq 1$. 

\newpage
\subsection*{Question 3}

\textbf{Prove:} \\

For $a \in X \setminus \bigcup \mathcal{F}$, we have $a \in X$ and $a \notin \bigcup \mathcal{F}$, by the definition of set difference.\\
Since $a \notin \bigcup \mathcal{F}$, $a$ is not in any set $Y \in \mathcal{F}$ by the definition of union.\\
Thus, for every $Y \in \mathcal{F}$, we have $a \in X \setminus Y$.\\
This implies $a \in \bigcap \{ X \setminus Y \mid Y \in \mathcal{F} \}$, by the definition of set difference.\\
Since $a \in X \setminus \bigcup \mathcal{F}$, by the definition of intersection,
\begin{align*}
X \setminus \left( \bigcup \mathcal{F} \right) &\subseteq \bigcap \{ X \setminus Y \mid Y \in \mathcal{F} \}.
\end{align*}

For $b \in \bigcap \{ X \setminus Y \mid Y \in \mathcal{F} \}$, for every $Y \in \mathcal{F}$, we have $b \in X \setminus Y$ by the definition of intersection. Which means, $b \in X$ and $b \notin Y$ for every $Y \in \mathcal{F}$.\\
Since $b$ is not in any $Y \in \mathcal{F}$, $b \notin \bigcup \mathcal{F}$, by the definition of union.\\
Since $b \in X$ and $b \notin \bigcup \mathcal{F}$, we have $b \in X \setminus \bigcup \mathcal{F}$.\\
 Because $b \in \bigcap \{ X \setminus Y \mid Y \in \mathcal{F} \}$, by the definition of set difference,
\begin{align*}
\bigcap \{ X \setminus Y \mid Y \in \mathcal{F} \} &\subseteq X \setminus \left( \bigcup \mathcal{F} \right).
\end{align*}

From both inclusions, we can conclude that:
\[
X \setminus \left( \bigcup \mathcal{F} \right) = \bigcap \{ X \setminus Y \mid Y \in \mathcal{F} \}
\]
\newpage
\subsection*{Question 4}
\begin{enumerate}
\item Reflexivity\\
Reflexivity is that for every $A \in \mathcal{P}(X)$,  $(A, A) \in R$.\\
Since the power set $\mathcal{P}(X)$ contains the empty set $\emptyset$ as an element, when $A = \emptyset$, we have $\emptyset \cap \emptyset = \emptyset$. This does not satisfy the condition $A \cap A \neq \emptyset$.\\
Therefore, $(\emptyset, \emptyset) \notin R$.\\
The relation $R$ is not reflexive.

\item Symmetry\\
Symmetry is that for every $A, B \in \mathcal{P}(X)$,  there is also $(B, A) \in R$.\\
 Assume $(A, B) \in R$, which means $A \cap B \neq \emptyset$.
 Since that $A \cap B = B \cap A$, if $A \cap B \neq \emptyset$, then$B \cap A \neq \emptyset$ is also ture. By the definition of $R$, the condition $B \cap A \neq \emptyset$ implies that $(B, A) \in R$.
\\The relation $R$ is symmetric.

\item Antisymmetry\\
Antisymmetry is that for every $A, B \in \mathcal{P}(X)$, whenever $(A, B) \in R$ and $(B, A) \in R$, it follows that $A = B$.\\
We can proof that $R$ is not antisymmetry by giving a counterexample of this.\\
Let $X = \{0,1, 2\}$(which is not empty), $A = \{0,1\}$ and $B = \{1, 2\}$.\\
    \[ A \cap B =   \{0,1\}\cap\{1, 2\} = \{1\} \]
Since $A \cap B \neq \emptyset$,  $(A, B) \in R$.
    \[ B \cap A =\{1, 2\}  \cap\{0, 1\}  = \{1\} \]
Thus, $B \cap A \neq \emptyset$,  $(B, A) \in R$.
However, $A \neq B$.
So, according to the definition, the relation $R$ is not antisymmetric.

\item Transitivity\\
Transitivity is that for every $A, B, C \in \mathcal{P}(X)$, whenever $(A, B) \in R$ and $(B, C) \in R$, it follows that $(A, C) \in R$.
We can proof that $R$ is not transitivity by giving a counterexample of this.\\
Let $X = \{1, 2, 3, 4\}$, $A = \{1, 2\}$, $B = \{2, 3\}$, and $C = \{3, 4\}$.
    \[ A \cap B = \{1, 2\} \cap \{2, 3\} = \{2\} \neq \emptyset \]
   So, $(A, B) \in R$.
    \[ B \cap C = \{2, 3\} \cap \{3, 4\} = \{3\} \neq \emptyset \]
    So, $(B, C) \in R$.
    \[ A \cap C = \{1, 2\} \cap \{3, 4\} = \emptyset \]
Since the intersection is the empty set, $(A, C) \notin R$.
So, according to the definition, The relation is not transitive.
\end{enumerate}

\newpage
\subsection*{Question 5}
(a) \\
Let\\
The probability that an attack event occurs: $P(A) = 0.02$\\
The probability that no attack event occurs:$P(\overline{A}) = 1 - P(A) = 0.98$\\
The probability that an event that the IDS raises a high-priority alert occurs:$P(R)$\\
The probability that an event that the IDS does not raise a high-priority alert occurs:$P(\overline{R})$\\
Using Bayes' Theorem:\\
\[
P(A \mid \overline{R}) = \frac{P(\overline{R} \mid A) \cdot P(A)}{P(\overline{R})}
\]
Since that
\[
P(\overline{R} \mid A) = 1 - P(R \mid A) = 0.02 
\]
\begin{align*}
P(\overline{R}) &= P(\overline{R} \mid A)P(A) + P(\overline{R} \mid \overline{A})P(\overline{A}) \\
&= 0.02\times0.02 + 0.95\times0.98 = 0.9314
\end{align*}
Thus,
\[
P(A \mid \overline{R}) = \frac{0.02 \times 0.02}{0.9314} = \frac{0.0004}{0.9314} \approx 0.00042946
\]

(b)\\
Using Bayes' Theorem:\\
\[
P(A \mid R) = \frac{P(R \mid A) \cdot P(A)}{P(R)}
\]
\begin{align*}
P(R) &= P(R \mid A)P(A) + P(R \mid \overline{A})P(\overline{A}) \\
&= 0.98\times0.02 + 0.05\times0.98 = 0.0686
\end{align*}
Thus,
\[
P(A \mid R) = \frac{0.98 \times 0.02}{0.0686} = \frac{0.0196}{0.0686} \approx 0.28571429
\]

(c)\\
$X$ is the number of days (out of 30) that the IDS raises an alert. Since days are independent, and there is two possible outcome for this (IDS raises a high-priority alert or not), this is a Bernoulli trial.\\

The expected value of a random variable $X$ on the sample space S is defined as:
\[
E(X) = \sum_{s \in S} p(s)X(S)
\]
For a binomial random variable $X$ , this simplifies to:
\[
E(X) = n \cdot p
\]
For $X$,$n = 30, p = 0.0686$, so
\[
E(X) = 30 \times 0.0686 = 2.058
\]
The variance of a random variable $X$ using the Theorem 
\[Var(X) = E(X^2) - E(X)^2
\]
For a binomial random variable:
\[
Var(X) = n \cdot p \cdot (1 - p)
\]
Substituting the values:
\[
Var(X) = 30 \times 0.0686 \times (1 - 0.0686) \approx 1.917
\]
(d)\\

\[
E[Y] =E[\text{Damage Cost}]  +E[\text{Responding Cost}]
\]


\begin{align*}
E[\text{Damage Cost}] &= 150000 \times P(A) = 150000 \times 0.02 = 3000 \\
E[\text{Responding Cost}] &= 5000 \times P(R) = 5000 \times 0.0686 = 343
\end{align*}

Thus,
\[
E[Y] = 3000 + 343 = 3343
\]


\end{document}